\subsection{内部対話 $A_d$ の分離：$A_{d,\mathrm{ego}}$ と $A_{d,\mathrm{wisdom}}$}
従来のNLPでは内部対話を一括して $A_d$（Auditory Digital）と表記してきたが、神経系疑似ベイズモデルにおいては、それが「どのネットワークで収束した確信 $\Omega$ を言語野が後付け翻訳したものか」によって厳密に二分される。

\begin{itemize}
  \item \textbf{$A_{d,\mathrm{ego}}$（エゴの言語化）}：\\
  DMNが確定させた $\Omega_{\mathrm{ego}}$ を、左脳のインタープリター（言語野）が後付けで言語化した独白。「あいつは私を軽視している」「どうせ失敗するに決まっている」といった、脅威と利害の反芻ループを再強化する内的言語である。
  \item \textbf{$A_{d,\mathrm{wisdom}}$（叡智の言語化）}：\\
  CEN主導で構造化された $\Omega_{\mathrm{wisdom}}$ を、論理的・客観的命題として言語化した内的思考。「このプロセスのボトルネックはここだ」「状況を別の視点から捉え直そう」という、現実的解決を推進する内的言語である。
\end{itemize}



\subsection{疑似ベイズ推論機の過剰強化：ループの閉塞と歪曲・幻覚の創発}

従来のNLPにおける「ストラテジー（Strategy）」とは、本来、特定の目標や行動を達成するために感覚表象を一定の順序で呼び出すフィードバック・アルゴリズムである。
TOTEモデル（Test-Operate-Test-Exit）に代表されるように、健全なストラテジーは「Exit（終了条件）」を満たすことでループを脱出し、次のタスクへとリソースを解放する。

しかし、このストラテジーを神経系疑似ベイズ推論機の逐次更新として捉え直したとき、工学的・病理学的な重大な脆弱性が浮かび上がる。
すなわち、\textbf{「適切なExit条件を欠いたストラテジーは、自己強化型の閉ループ（ポジティブ・フィードバック）へと縮退し、事前予期を異常に過剰強化（Over-tuning）させる」}という点である。

\subsubsection{自己参照ループによる事前予期 $\hat{\Omega}$ の極小過学習}
逐次ベイズ更新の方程式を想起されたい。
\begin{equation}
    P(H_k \mid D_1, \dots, D_m) \propto P(H_k) \prod_{j=1}^m P(D_j \mid H_k)
\end{equation}
外部環境からの客観的なテスト（反証データ）によってループが遮断されない場合、推論機は直前の事後確率をそのまま次の事前確率へと無制限に繰り越していく。
エゴの脅威仮説 $\hat{\Omega}_{\mathrm{threat}}$ に基づくループがExitを持たずに回転し続けると、事前確率は極限まで肥大化し、仮説空間の自由度は失われる。

\begin{equation}
    P(\hat{\Omega}_{\mathrm{threat}}) \longrightarrow 1
\end{equation}

事前確率がここまで極端に偏ると、尤度計算において致命的な「歪曲（Distortion）」が発生する。
本来はその仮説と無関係であるはずの微小な環境ノイズやランダムな刺激であっても、推論機はそれを強引に自らの仮説に合致する強力なエビデンスとして解釈し、過剰マッチング（過学習）を起こすのである。

\subsubsection{自臭症のメカニズム：些細なノイズから幻覚への連続体}
この過剰強化の典型例として、臨床的に「自臭症（自己臭恐怖症）」と呼ばれる病態の神経ダイナミクスを挙げることができる。

自臭症の患者においては、「自分は周囲に不快な悪臭を放っているのではないか」という自己保全的脅威ナラティブ $\hat{\Omega}_{\mathrm{odor}}$ が、Exitのないループとして内部で回転している。
\begin{enumerate}[label=(\arabic*)]
  \item \textbf{Exitの欠落}：嗅覚は急速に疲労・順応するため、本人は自らの体臭を客観的に検証（Exit判定）できない。検証不能のまま推論が空転し続ける。
  \item \textbf{事前予期の肥大化}：反証が得られないため、逐次更新によって「自分は臭っている」という事前確率 $P(\hat{\Omega}_{\mathrm{odor}})$ が極大化する。
  \item \textbf{知覚ノイズの強制歪曲}：周囲の人間が「咳払いをした」「鼻を触った」「窓を開けた」といった無関係な視覚・聴覚ノイズ $v, a$ が、すべて「悪臭に対する忌避行動」という極めて高い尤度を持つエビデンスへと歪曲される。
\end{enumerate}

さらにこの過剰強化が臨界点を超えると、単なる解釈の歪曲（関係念慮）にとどまらず、感覚知覚の生成そのものに破綻が生じる。
電車内の遠くの話し声や空調の機械音といった、意味を持たないランダムな聴覚ノイズ $a_{\mathrm{noise}}$ に対し、肥大化した事前予期が強力なトップダウン信号として感覚皮質へ直接書き込まれる。
歪曲が生じるこのプロセスを$:\Longrightarrow$で表す。
\begin{equation}
    a_{\mathrm{noise}} :\xrightarrow[\hat{\Omega}_{\mathrm{odor}}]{\text{Over-tuned Prior}} \Omega \text{ (自分は臭い) } :\Longrightarrow A_e \text{（臭い）}
\end{equation}

物理的には単なる空気の摩擦音や雑音にすぎないものが、神経系内部では明瞭な「他者の悪口」という外的表象 $A_e$ として確定・知覚されてしまう。
これは、推論のトップダウン予測がボトムアップの感覚入力を完全に圧倒・上書きした状態であり、\textbf{統合失調症などに見られる「幻聴（Auditory Hallucination）」の直接的な萌芽}にほかならない。

幻聴や妄想は、何もない無の空間から突然発生するオカルト現象ではない。
Exitのない自己強化ループによって疑似ベイズ推論機の事前予期が暴走し、日常の些細な環境ノイズを強引に特定の破滅的ナラティブへ過剰マッチングさせた結果として生じる、計算論的神経システムの構造的バグなのである。


\subsection{現実見当識 $\Omega_{\mathrm{neutral}}$ の侵食と精神病理への縮退}

さらに深刻な病理的力学は、これら激しく対立・振動するエゴや信念のナラティブ確信（$\Omega_1, \Omega_2$ など）が、自己と世界の物理的基盤を担保している「現実見当識の確信（$\Omega_{\mathrm{neutral}}$）」そのものと正面から矛盾・衝突を起こした際に発現する。

通常、人間が日常を維持できているのは、物理空間の常識や身体の限界を規定する $\Omega_{\mathrm{neutral}}$ に対する事前予期 $\hat{\Omega}_{\mathrm{neutral}}$ が、経験的な逐次更新によって極めて高く維持されているからである。
しかし、外部からの知覚エビデンス $\{(VAKOG)_e\}$ や、情動的身体反応 $\Lambda$ の強烈なループバックが、現実見当識とは相容れない強固なナラティブ仮説――例えば「私は見えない超常的存在に監視・操作されている」「この部屋の空間自体が私を呪詛している」といった極端なマクロ仮説――を過剰に支持し続けた場合、神経系疑似ベイズ更新の力学系は致命的な転換点を迎える。

排他的な仮説空間において、ある極端なナラティブの事後確率が不当に跳ね上がれば、ベイズの正規化制約により、対立する基盤的確信の確信度は必然的に押し下げられる。

この逐次更新の破綻プロセスが反復されると、現実見当識を支えていた事前予期 $\hat{\Omega}_{\mathrm{neutral}}$ 自体が不可逆的に減衰し、最終的に生体は客観的現実のパースペクティブを完全に喪失するに至る。
本稿では、これを神経系疑似ベイズ推論機における\textbf{「現実見当識の確率的縮退モデル」}として、精神病（Psychosis）や精神崩壊の一つの発症機序として仮説づける。

もちろん、精神疾患の全容がこの単一のモデルのみで説明できるわけではない。
しかし、この破綻現象は、極めて強固で絶対的なドグマを内面化した宗教的信念の信奉者や、カルト的共同体の成員において顕著に観察される。
外的現実と真っ向から対立する超常的教義が、強烈な儀式や情動体験（$\Lambda$ の過剰励起）を通じて絶対的な確信 $\Omega$ としてロックインされたとき、物理的な現実世界のエビデンスはすべて「悪魔の誘惑」や「神の試練」として歪曲解釈（過学習による尤度歪曲）される。
その結果、現実見当識を司る $\Omega_{\mathrm{neutral}}$ は完全に確率を失って破綻し、客観的対話が一切成立しない重篤な妄想体系・現実遊離状態へと滑落していくのである。


